% SOURCE docs/release/2026-10-29/source/docs/writing-repair/2026-10-29/src/1ad8b0cde9c2e6ff.tex CHAR 0
\chapter{Translating a banking decision into a causal question}
\label{la:translation}

The preceding evidence review leaves a practical starting point: assistance
must preserve the question before it can improve an answer. This chapter
builds that question from population, action, outcome, loss and information
time. Incentive and servicing examples expose common substitutions; a
literature-graph study separates extracted claims from hypotheses; the
formulation workflow then leads to robust population comparisons and goal
clarification. Together these treatments explain which commitments an analyst
must resolve and which alternatives an assistant must retain.

\section{A business sentence can contain several questions}

``Should we increase the deposit offer?'' sounds like a single question, but
it leaves several choices unresolved. The intended action could be a change
in the quoted rate, the eligibility rule, the message shown to a customer or
the allocation of a fixed promotional budget. The outcome could be a new
deposit, retained balances, funding cost or total relationship value. The
population could be existing customers, applicants who received an offer or
everyone eligible under a proposed rule. Each choice defines a different
causal problem.

A language assistant can expose these alternatives and organize the evidence
needed to resolve them. It should not choose the easiest target to estimate
and silently treat that target as the owner's decision. In the project funding
case, the response of balances to a 25-basis-point increase and the change in
net funding cost over thirty days are retained as separate candidates. Their
business motivation is shared; their outcomes and decision consequences are
not identical. The unresolved choice is useful information, not a failure
that must be hidden by returning a number.

The co-pilot roadmap discussed by Alaa and colleagues motivates this progression
from a scientific question through assumptions, statistical targets and
interpretation. Their illustrative medical examples attend to eligibility,
baseline time, treatment and follow-up. The inspected paper also reports that
generic uncued answers can omit relevant design concerns. Applying those
principles to banking is a project proposal; the medical examples do not
validate banking clocks or interventions. \cite{v12_la_2407_19118}

\section{A question as a collection of commitments}

The project represents a question as
\begin{equation}
 Q=(\mathcal U,P,\mathcal A,\pi,\pi_0,Y,H,\nu,\ell,\mathcal I).
 \label{la:question}
\end{equation}
Here \(\mathcal U\) is the unit of analysis; \(P\) is the target population;
\(\mathcal A\) is the feasible action set; \(\pi\) and \(\pi_0\) are the
policy and comparator; \(Y\) is the outcome; \(H\) is the horizon; \(\nu\)
is the measure used to average over the population; \(\ell\) is the loss,
including its units; and \(\mathcal I\) specifies information available at
each action. This notation defines the project's declaration contract. It is
not an identification theorem.

The unit needs particular care. An account can be an observation while a
customer is the decision unit and a household is the unit within which actions
interact. An assistant can ask which role each identifier plays, but a column
called customer ID does not certify ownership or stable linkage. If several
accounts share a budget or pricing rule, an account-level action may also be
an infeasible description of the intervention. The feasible action set should
reflect what the responsible institution can actually change.

The comparator must be as explicit as the intervention. ``No treatment'' might
mean no new contact, continuation of the current contact policy, no rate
increase, or a different offer under a fixed budget. A dynamic policy can
include future actions conditional on history, so its comparator is also a
history-dependent rule. An effect relative to one comparator cannot be
reused for another solely because the same treatment variable appears in the
dataset.

The outcome and loss should also remain separate. A change in balances has
units of currency held over a horizon; funding expense combines balances,
rates and time. A larger deposit balance can have a positive acquisition
interpretation and an unfavorable net-cost implication. If \(g(a)\) is a
specified gross benefit and \(c(a)\) is a specified cost, the decision may
depend on \(E[g(a)-c(a)]\), not on \(E[g(a)]\) alone. The definitions and
valuation assumptions must be supplied before that subtraction is meaningful.
The assistant should make missing costs visible rather than invent prices.

\section{Changing a population measure changes the answer}

Let \(X\) denote baseline features and let
\(\tau(x)=E[Y(1)-Y(0)\mid X=x]\). A weighted target can be written
\begin{equation}
 \theta_w=\frac{E[w(X)\tau(X)]}{E[w(X)]},\qquad E[w(X)]>0.
 \label{la:weighted-target}
\end{equation}
Uniform, treated-population, overlap and trimmed targets use different
weights. Under the conditions that make a propensity \(e(x)=\Pr(A=1\mid X=x)\)
appropriate for the intended population, familiar choices are \(1\),
\(e(x)\), \(e(x)(1-e(x))\) and
\(1\{a\leq e(x)\leq b\}\), respectively. The formula records which average
is wanted. It does not establish exchangeability, support, estimation rates
or coverage.

An exact project example makes the distinction visible. Let two strata have
equal population probability, with propensities \(0.2\) and \(0.5\), and
effects \(0\) and \(10\). Trimming to \([0.1,0.9]\) keeps both strata, so
the trimmed-population effect is \(5\). Overlap weights are \(0.16\) and
\(0.25\), giving
\[
 \theta_{\mathrm{overlap}}=
 \frac{0.16(0)+0.25(10)}{0.16+0.25}=\frac{250}{41}
 =6.09756\ldots.
\]
Both calculations can be exact and still answer different questions. A label
such as ``overlap population'' is insufficient if its underlying weighting
rule is not retained.

This issue is directly relevant to the inspected OCI agent paper. Chou and
colleagues describe a principal/actor/critic workflow, reproducible notebooks,
point-in-time estimation and overlap diagnostics. The inspected trimming
passage uses a target label that must be reconciled with the actual weighting
rule before mapping it to conventional overlap weighting. The proposed
architecture is informative, but an assistant should not inherit a target
label without checking its definition. \cite{v12_la_2607_22443}

Changing the target to remove unsupported accounts can be scientifically
reasonable, but it is a new decision question. A bank owner may care about
precisely those accounts for which treatment support is weak. An honest report
can say that a narrower target is estimable under the supplied assumptions
while the original target remains unresolved. It should show the excluded
population, the new averaging rule and the implications for the decision.

\section{Information is governed by an availability clock}

The date on which an event occurred differs from the date on which an analyst
or policy could observe it. For a recorded variable \(V\), write \(e_V\) for
event time and \(a_V\) for availability time. The finite project fixtures
require \(e_V\leq a_V\), and a policy acting at time \(t\) can use the value
only if \(a_V\leq t\). A derived feature becomes available no earlier than
its required inputs. These rules prevent an apparently historical feature
from leaking future information into a decision.

In the servicing fixture, an account condition occurs at day ten but becomes
available at day eleven. A rule acting on day eleven may use it; a rule acting
on day ten may not. The distinction matters even when both values later appear
in the same warehouse table. A language assistant should request the recording
and availability convention, rather than infer it from a field name or a
sorted timestamp.

Advance announcements and forecasts require different variable definitions.
A forecast observed today of next month's condition is a current forecast
variable, not a measurement of a future realized condition with an impossible
availability time. Revised historical data similarly require a vintage or
revision convention. The project's finite validator does not implement every
revision protocol. It checks the stated recorded-value contract and leaves
more complex clocks as explicit extensions.

Equal timestamps do not determine direction. Two columns aggregated to the
same day may hide a finer action sequence or a common cause. The assistant
can propose a finer clock and identify the evidence needed to support it.
It cannot establish temporal precedence by placing one variable first in a
diagram. When the granularity is insufficient, the ambiguity should survive
into the model class.

\section{An available covariate is not automatically an adjustment variable}

In a sequential servicing model, \(A_0\) changes an intermediate condition
\(L_1\), which affects both a later action \(A_1\) and the outcome \(Y\).
The same \(L_1\) is a mediator for the first action and a common cause of the
second action--outcome relation. Its appropriate use depends on the target.
Including every available column in one regression does not preserve all
possible total, direct and policy effects.

The project declaration therefore distinguishes a policy-history purpose
from an ordinary adjustment purpose. A simple directed-path diagnostic flags
a post-treatment descendant proposed as an ordinary total-effect control.
It does not declare all post-treatment conditioning invalid. Controlled
direct effects and longitudinal procedures need their own question and
formal analysis. A warning from a declaration validator is not a replacement
for an adjustment criterion or an identification algorithm.

Observation processes can alter the target in a similar way. Suppose only
successfully linked records enter an analysis and both action and repayment
affect successful linkage. The selection indicator is part of the causal
problem, and an effect among linked records may differ from an effect in the
original population. A record-link score is an observed proxy; it is not
verified ownership. The observation fixture deliberately retains a latent
ownership variable, a selection node and an unresolved population choice.

CausalDS provides a useful literature distinction here: its conceptual SCM
variables and its separate observation model are not interchangeable. The
paper's narrated synthetic scenes can test how an agent handles that
distinction, but they do not supply authentic banking semantic or assignment
truth. \cite{v12_la_2607_08093}

\section{The handoff to formal analysis}

The assistant's useful output is a reviewable declaration and a record of
what remains unresolved. A selected declaration contains one explicit target,
its population measure, the intervention and comparator, clocks, feasible
actions, outcome and loss. An ambiguous declaration retains candidate targets
without covertly selecting one. A contradicted assumption cannot be left as
an inconsequential note while the calculation proceeds as if it held.

The local project validator exercises these distinctions with four synthetic
banking cases and adverse contract tests. It checks declared identities,
types, references and selected clock/graph conditions. It does not authenticate
an institution, certify a reviewer, determine whether a causal assumption is
true, or implement every MAG/PAG/CPDAG validity condition. A declaration ready
for formal analysis is ready for the next conditional operation, not a
qualified bank recommendation.

This boundary also explains the limited scope of the executable operation
tested with these declaration fixtures. A declaration can describe a dynamic
regime while that selected operation supports only a static binary population
ATE. The correct handoff then reports unsupported target mapping. Later
fitted-summary and servicing studies have their own declared regimes; their
existence does not qualify this particular mapping. Translating that request into a
static regression would produce a result at the cost of changing the user's
question. The appropriate remedy is to qualify a longitudinal operation or
obtain an explicit new target decision, with the original question preserved.


% SOURCE docs/release/2026-10-29/source/docs/writing-repair/2026-10-29/src/4bcfb300b6807e43.tex CHAR 0
\section{A complete translation example: the incentive question}
\label{v13:la-dose}

Suppose a manager asks an agent to ``find the incentive uplift and choose the
best offer.'' The sentence leaves at least three mathematical objects open:
a binary offer effect, a continuous local slope, and a finite-dose policy value.
The agent can help by presenting these alternatives with their populations,
units and action constraints. It should not select whichever object happens
to be exposed by the first available package.

Let $A$ be incentive intensity in a declared unit and let $Y$ be retained net
value over a fixed horizon. The finite question is
$\Delta(a_1,a_0)=\mathbb E\{Y(a_1)-Y(a_0)\}$. A local derivative, where it
exists, is $\mu'(a)$. A covariance slope averages features of the response
under an exposure distribution and need not equal either object. The inspected
grf and dose-response interfaces make this a practical interoperability issue,
not merely a linguistic distinction. \cite{v13_sw_grf,v13_sw_causaldrf}

For the stipulated curve $\mu(a)=20a-3a^2$, a uniform exposure on $[0,2]$
gives covariance slope 14, while moving from dose one to dose two changes the
mean by 11 CNY. An agent that labels 14 as the value of the specified change
has mistranslated the problem even if it calls its numerical tool correctly.
The exact project calculation makes the error detectable without a live model
experiment. \cite{v13_project_exact}

A useful translation record would therefore state: eligible existing customers;
the information available before the offer; current dose one and feasible dose
two; a fixed retention horizon; treatment and outcome currency units; and
whether incentive expense is already included in $Y$. If expense is excluded,
the decision value is $\Delta(a_1,a_0)-\{c(a_1)-c(a_0)\}$. Subtracting it again
when $Y$ already measures net value would double-count cost. The record also
needs account closure/default definitions and a sampling unit for uncertainty.

The agent's role can be useful at each handoff. It can flag unspecified units,
find source definitions, propose alternative causal questions and generate
typed requests. Formal checks can verify that a numerical result matches the
declared query. Domain owners must still confirm the feasible intervention and
measurement meanings; neither fluent prose nor a schema establishes them.

\section{A servicing policy is a rule with an information boundary}

For sequential servicing, an action rule uses $H_t$, the information actually
available at decision time. A data field recorded later may describe an earlier
event while remaining unavailable to the operating policy. An agent assembling
history features must preserve both event and availability times. Otherwise
the proposed policy uses information the bank would not possess when acting.

The causal query should distinguish evaluating a specified rule, choosing among
rules and claiming an optimal hidden-state policy. A longitudinal simulator can
evaluate supplied regimes under its assumptions. Its ability to accept a custom
history function does not identify private customer states or make an unknown
logger known. The inspected pygformula class exposes useful histories, event
and regime options, but its source interface alone establishes neither policy
optimality nor uncertainty under arbitrary dependence. \cite{v13_sw_pygformula}

A concrete policy-selection hazard can be shown without an asymptotic argument.
Suppose two equally valuable policies each have independent evaluation error
equal to minus one or plus one with equal probability. The maximum of the two
errors is minus one only when both are minus one, and plus one otherwise.
Its expectation is therefore $\tfrac34-\tfrac14=\tfrac12$. Reporting the
selected policy's same-data evaluation is optimistic. If the selected policy
is evaluated with a fresh independent error of mean zero, the evaluation error
has mean zero conditional on the selection. Exact enumeration verifies both
claims. \cite{v13_project_exact}

This example does not prove that a particular split is sufficient for a
dependent servicing panel. It explains why the selection and evaluation laws
must be specified. An agent can preserve the split, freeze candidate rules and
refuse reuse of a held-out set for tuning. A scientific argument must still
justify independence or the required dependence control. Training a policy
and then narrating its favorable development score is not an independent
evaluation.

\section{Translation outputs that deserve abstention}

Some business questions should produce a request for more information or a set
of alternatives. If ``the customer'' may mean the legal entity, the beneficial
owner or a linked household, choosing one silently changes both the population
and interference structure. If the intended action is not operationally
feasible, the mathematically optimal action is irrelevant. If the loss function
and costs are unknown, the agent can compare statistical targets but cannot
assert an investment recommendation.

Abstention should be specific. It should identify the missing definition,
record the alternative interpretations and explain which downstream operation
depends on the answer. That is different from a generic refusal to analyze.
For an incentive query, the missing item may be treatment units or whether
cost is already included. For servicing it may be the information clock or
the policy class. For interference it may be network direction or the allocation
policy. These are concrete commitments that owners can resolve and numerical
tools can then preserve.


% SOURCE docs/release/2026-10-29/source/docs/writing-repair/2026-10-29/src/7d702eb060b9bf4a.tex CHAR 0
\section{From a literature graph to a testable banking hypothesis}
\label{exec20:hypothesis-scope}
A bank translating a business question into a causal analysis needs more than a plausible arrow. Suppose a servicing team asks whether a digital reminder reduces missed payments among customers whose income has become irregular. Literature can suggest mechanisms involving attention, financial stress, access and payment timing. A useful assistant should expose those possibilities and their sources, then help specify an intervention, eligible population, outcome horizon and decision clock. It should not quietly turn a sentence about association into an identified effect of sending a reminder. Tong and colleagues provide a concrete design for using language models and a literature graph to generate research hypotheses. Their psychological setting is informative about the design of that assistance, while supplying no validation of the banking intervention.~\citep{exec20_hypothesis_generation}

\subsection{What the literature graph represents}
\label{exec20:hypothesis-claims}
The inspected July 2024 article describes three stages: retrieve literature, extract relation pairs and generate hypotheses from predicted links. From a wider psychological and neuroscience collection, the authors select 43,312 articles using a journal-title criterion; 35,797 come from \emph{Frontiers in Psychology}. That concentration is a property of the selected evidence, not a census of psychological knowledge. PDF cleaning and chunks of at most 4,000 tokens precede GPT-4 extraction. The prompt explicitly requests both causal and correlational relations, their concepts, and the sign of causal relations. Model filtering and periodic expert feedback follow. The resulting graph is reported to contain about 197,000 concepts and 235,000 connections. Neo4j supplies storage; node2vec embeddings and neighbor Jaccard similarity guide candidate-link selection. GPT-4 then turns selected concept pairs into hypotheses. These are the authors' reported components and counts, not a reproduced implementation.~\citep{exec20_hypothesis_generation}

A bank-facing evidence object should retain more information than the concept pair. Let $q$ identify a source claim, $d_q$ its document and version, $s_q$ the inspected passage, $r_q$ its asserted relation type, and $c_q$ its population, timing and study-design context. A useful proposed representation is
\begin{equation}
  e_q=(d_q,s_q,r_q,c_q,\hbox{extraction status},\hbox{review status}).
  \label{exec20:hypothesis-evidence}
\end{equation}
The claim graph constructed from these objects need not be an acyclic structural causal model. Different studies can use different populations, variable definitions and time resolutions; an article can report a causal interpretation whose identifying assumptions remain disputed. Deduplication must preserve disagreement and source independence. A collection of paraphrases of one study cannot acquire the status of several replications. Model confidence and an expert's agreement that the text contains a claim are also different quantities from evidence that the claim is true in the proposed bank population.

The paper's extraction check makes this distinction measurable. Among 238 pairs classified as causal by GPT-4, the main article reports 156 agreed causal labels and 31 pairs not recognized as relationships. The complementary 51 therefore have a different recognized relation label under that classification. Consequently,
\begin{equation}
 \frac{156}{238}\simeq 65.55\%,\qquad
 \frac{207}{238}\simeq 86.97\%,\qquad
 \frac{31}{238}\simeq 13.03\%.
 \label{exec20:hypothesis-denominators}
\end{equation}
The small rounding differences from the article's displayed percentages do not change the substantive point: relation-existence agreement is a different endpoint from causal-label agreement. Neither denominator measures an intervention's ground-truth effect. The authors' familiar-paper evaluation and feedback process supplies useful exploratory evidence; it does not establish performance on previously unseen banking mechanisms.~\citep{exec20_hypothesis_generation,exec20_hypothesis_reference}

\subsection{Why a predicted link remains a hypothesis}
\label{exec20:hypothesis-graph}
For neighbor sets $N(a)$ and $N(b)$, the usual Jaccard score is
\begin{equation}
 J(a,b)=\frac{|N(a)\cap N(b)|}{|N(a)\cup N(b)|}.
 \label{exec20:hypothesis-jaccard}
\end{equation}
This expresses shared graph context under a chosen neighbor convention. A denominator of zero needs an explicit implementation policy. Direction, edge types, duplicate claims and temporal aggregation also affect what a neighborhood means. A score in $[0,1]$ is not thereby a probability of causation, a confidence level or an effect magnitude.

An authored counterexample illustrates the missing implication. Take mutually independent zero-mean $Z,\varepsilon_A,\varepsilon_B,\varepsilon_P$ and the structural equations
\begin{equation}
 A=Z+\varepsilon_A,\qquad B=Z+\varepsilon_B,
 \qquad P=A+B+\varepsilon_P.
 \label{exec20:hypothesis-world}
\end{equation}
The DAG has $Z\to A$, $Z\to B$, $A\to P$ and $B\to P$. Using undirected adjacency neighborhoods, $N(A)=N(B)=\{Z,P\}$, so $J(A,B)=1$. Yet intervening on $A$ changes neither $Z$ nor the equation for $B$: the derivative of $E[B\mid\operatorname{do}(A=a)]$ with respect to $a$ is zero. Shared cause and shared consequence create identical neighborhoods without a direct or total effect of $A$ on $B$. This is a statement about the displayed construction; the uninspected package's neighbor convention is not assumed.~\citep{exec20_hypothesis_reference}

For a banking interpretation, let $A$ encode a reminder policy, $B$ a repayment response, and $Z$ a common predisposition in a hypothetical observational population. If a planner wrongly treats similarity one as a success probability for an intervention worth USD 100 at cost USD 20, the apparent margin is USD 80. Under the stipulated zero-effect law, the incremental margin is USD -20. No bank profit or package failure is measured here. The construction shows why a candidate generator must return a proposed relation with its evidence, while a subsequent identification step supplies any admissible causal effect or refusal.

\subsection{What the reported evaluation can establish}
\label{exec20:hypothesis-evaluation}
The authors focus on well-being and generate 130 hypotheses. The principal comparison uses 30 randomly selected model-and-graph hypotheses, 30 expert-selected hypotheses, 30 hypotheses sampled from 41 ideas supplied by 13 responding doctoral students, and 30 Claude-2 hypotheses selected by GPT-4 from 50 generated ideas. Three experts rate novelty and usefulness, with hypotheses presented in randomized order. These pipelines differ in generation and selection. The comparison therefore concerns the resulting pipelines and selected ideas, rather than isolating one graph component under identical resources. Reported reviewer correlations vary, including a novelty correlation of 0.069 for one reviewer pair; an average rating cannot erase that disagreement.~\citep{exec20_hypothesis_generation}

The paper reports higher novelty for its randomly selected graph-assisted ideas than for the Claude comparison. The expert-selected arm is not uniformly significant: Table 5 gives mean-comparison $p=0.085$ and median-comparison $p=0.550$ against Claude. A later 60-versus-60 GPT-4 ablation reports a novelty advantage but a usefulness comparison of $p=0.1937$. Nonsignificance does not prove equal usefulness, nor does an insignificant human-versus-model comparison establish equivalence to experts without a justified equivalence margin and suitable design. These endpoints are judgments of ideas; the proposed mechanisms were not validated by the reported rating experiment. BERT representations and t-SNE plots describe semantic structure and require their own interpretation. They do not replace an intervention experiment.~\citep{exec20_hypothesis_generation}

Several printed numerical statements also require reconciliation before the proposal inherits their inferential force. The usefulness paragraph pairs $F(3,116)=5.25$ with $p=0.553$. Independent evaluation of the ordinary $F$ upper tail gives approximately $0.001967$. In the semantic-distance analysis, the printed $F(3,1652)=84.1611$ is paired with explained variance 86.96\%. Under ordinary one-factor ANOVA with the printed degrees of freedom,
\begin{equation}
 R^2=\frac{d_1F}{d_1F+d_2}
     =\frac{3(84.1611)}{3(84.1611)+1652}
     \simeq 0.132573.
 \label{exec20:hypothesis-anova}
\end{equation}
These are arithmetic checks conditional on the stated conventional model, not a reconstruction of unavailable raw ratings. The correct repair is to obtain the analysis and resolve which printed quantities or definitions need correction. The prose also varies from its table for a human-comparison $t$ value and calls a median comparison significant where the table gives $p=0.550$. Hypotheses contributed by the same person, repeated ratings and distances sharing the same hypothesis create dependence questions; their inferential treatment cannot be recovered from a significance label alone.~\citep{exec20_hypothesis_generation,exec20_hypothesis_reference}

\subsection{A public evaluation before a bank prototype trial}
\label{exec20:hypothesis-prototype}
The proposed library can adopt the separation of literature extraction, graph-assisted candidate generation and expert evaluation as a design hypothesis. Direct code adoption remains premature: the selected article does not itself qualify immutable implementation identities, exact model snapshots, full prompts and thresholds, dependency licenses, or complete corpus redistribution rights. Its data statement makes only part of the evidence available with supplementary material and offers additional data on request. The earlier preprint's appendices are a different version and cannot silently substitute for the revised supplement. The article's own limitations acknowledge extraction error, synonymous concepts, a restricted hypothesis sample and a shift in scarce expert work from generation to evaluation.~\citep{exec20_hypothesis_generation}

A first public experiment should freeze a licensed corpus and evaluate four linked but separate questions. Can extraction recover source-supported claims, relation type, direction, sign and uncertainty on held-out documents? Can an assistant translate a business description into an operational population, intervention and outcome without inventing assignment or observability? Does graph-assisted generation improve blind ratings of useful, testable hypotheses relative to retrieval-only, language-model-only and human-assisted baselines under matched budgets? Finally, do proposals survive identification and intervention checks in generated banking mechanisms with known truth and in suitable publicly documented experiments? Novelty and semantic diversity remain secondary to correctness and decision loss when an action is proposed.

Document families, shared studies and generated mechanism families should be separated across development and evaluation to reduce leakage. The public protocol should retain contradicted claims, extraction failures, unavailable answers and refusal. Human evaluation needs a declared sampling unit, a predeclared treatment of reviewer disagreement and dependence, and a held-out adjudication process. Actual calls, retries, tokens, expert minutes and component compute should be recorded. The comparison should price the full review workflow rather than assume that automatic hypothesis production removes expert work.

Across library, CLI and MCP, the proposed output should distinguish \texttt{extracted\_claim}, \texttt{hypothesis}, \texttt{reviewed\_assumption} and \texttt{identified\_query}. It should carry immutable source anchors, event-time definitions, unresolved conflicts and the checks still required for promotion. A client requesting an effect must not receive a link score under an effect field. These contracts are a concrete route to a useful public prototype; qualification of a banking intervention still requires evidence about that intervention.



% SOURCE docs/release/2026-10-29/source/docs/writing-repair/2026-10-29/src/62cabebeedc23664.tex CHAR 0
\hypertarget{a-formulation-workflow-that-preserves-unresolved-alternatives}{%
\section{A formulation workflow that preserves unresolved
alternatives}\label{a-formulation-workflow-that-preserves-unresolved-alternatives}}

\label{exec13:formulation}

A language model can make a business question more precise by exposing
the decisions that its wording leaves open. Its useful output is a set
of explicit candidate questions, with evidence for their variables and
clocks, followed by a reasoned choice of a supported analysis. Fluency
alone does not determine which candidate the bank intended. A formal
causal engine also cannot repair a mistake that was made before the
graph, intervention and target reached the engine.

Consider a manager who asks, ``Which customers should receive a
retention invitation next month, and how much value will it create?''
One interpretation is the effect of offering the invitation to every
currently eligible customer. Another is the effect among customers who
historically received it. A third is the value of a targeting rule under
a fixed contact budget. The first two are averages over different
populations. The third also requires a feasible policy class, the
information available when the rule acts, a comparator policy and a cost
definition. A model that estimates historical recipients' average effect
and calls it the expected value of the new targeting rule has changed
the question.

The formulation record should therefore contain the eligible population
and exclusions, the sampling unit, invitation and comparator, actual
receipt if distinct, the decision time, permissible information, outcome
horizon, monetary units and accounting treatment of contact cost. It
should distinguish invitation from successful contact and from
subsequent product use. Those variables may be causally related, but
they are different interventions. It should also retain alternatives
that remain unresolved: all eligible customers versus historical
recipients, a single action versus a sequence, and gross revenue versus
net value.

For a binary invitation \(A\), a baseline record \(W\) available at time
zero, and a fixed-horizon outcome \(Y\), an initial target may be

\[
 \psi=\mathbb E_{\nu}\{Y(1)-Y(0)\},
\]

where \(\nu\) names the eligible population. The adjustment expression

\[
 \psi=\int\{m_1(w)-m_0(w)\}\,d\nu(w),
 \qquad m_a(w)=\mathbb E[Y\mid A=a,W=w],
\]

requires consistency, conditional exchangeability, positivity, suitable
sampling and the declared interference restriction. The model may help
elicit and document those assumptions. It cannot infer their truth from
a plausible explanation or from a successful estimator call. In a
generated development example, the project can stipulate and inspect the
complete assignment law. In a later bank application, the assignment
mechanism and measurement meanings need accountable scientific
ownership.

\hypertarget{what-the-frontier-evidence-supports}{%
\subsection{What the frontier evidence
supports}\label{what-the-frontier-evidence-supports}}

The knowledge-assisted perspective of Kıcıman and colleagues motivates
graph drafts, critique and connections to formal tools. Its own
memorization probes and discussion of novel scenarios limit what
benchmark success establishes. The inspected original distinguishes
knowledge-based graph generation from conventional discovery from
observations. For this proposal, a directional suggestion is a
hypothesis with source provenance, not a discovered bank fact. A model
can suggest that financial distress confounds retention contact and
subsequent deposits; the analysis still needs to determine what was
measured, when staff observed it and how it affected assignment.
\citep{v12_la_2305_00050}

The medical co-pilot proposal develops a causal roadmap from a
scientific question through study design and interpretation. Its
examples of prevalent versus incident users are directly relevant to the
distinction between existing account holders and newly eligible
customers. That transfer is a project design proposal, not evidence that
a medical demonstration validates banking workflows. The paper's
evaluation discussion also makes clear that medical question-answering
ability and causal study-design quality are different evaluation
targets. Carefully cued illustrative responses do not establish reliable
autonomous formulation. \citep{v12_la_2407_19118}

CLadder provides formally generated causal questions and explicitly
separates levels of causal reasoning. Its results and limitations
support using a formal reference to evaluate whether verbalized
information was translated into the right query. Binary answer accuracy
on its selected query families does not measure effect calibration,
correct bank populations or decision loss. The paper also identifies the
language-to-formal-tool interface as a remaining problem. The proposed
baseline therefore checks the structured query before it scores a
numerical answer. \citep{v12_la_2312_04350}

The 2026 CausalDS paper adds narrated synthetic structural models, data
files, an observation layer and separate private scoring information. It
distinguishes the conceptual variables used for identification from
their noisy observed proxies. Its reported results separate symbolic
reasoning, estimation, interval quality, abstention and tool cost. This
is useful evidence for a multidimensional evaluation: correctly
extracting a story's graph does not guarantee that the observed data
identify or precisely estimate the stated effect. The paper's historical
model results are author-reported; they have not been reproduced in this
project. \citep{v12_la_2607_08093}

The 2026 observational causal-inference agent paper provides a concrete
principal, actor and critic workflow built around deterministic analysis
tools and reproducible notebooks. It reports favorable comparisons under
its evaluated settings, while its disguised observational-study examples
retain important overlap-related refusals. This supports investigating
constrained analytical workflows with explicit diagnostics. It does not
establish unmeasured-confounding control, longitudinal-policy
identification or valid conditional coverage after the critic selects
which estimates to report. Its point-in-time scope must survive transfer
to the banking proposal. \citep{v12_la_2607_22443}

Two additional primary studies sharpen the boundary between an effective
workflow and causal truth. Sng and colleagues study inconsistencies in
causal-edge ratings and compare retrieval augmentation with a
debate-and-arbiter design. Their retrieval experiment supplies a mocked
true causal graph. The authors explicitly state that reduced
inconsistency does not imply accurate causal prediction and call for
true-graph validation. The title's language about eliminating
hallucinations should therefore not become a guarantee in this product
proposal. Consistent agents can agree on an incorrect edge, and a
retrieval corpus can encode the same incorrect premise.
\citep{exec13_hallucination}

CausalPulse describes query planning, semantic descriptions, process
constraints, agent cards and MCP/A2A coordination for manufacturing
diagnostics. Its reported operational success is a useful architecture
comparison. The paper separately reports root-cause attribution, with a
PSE Hits@1 of 0.44 and F1 of 0.39, using annotations of variables that
changed at the anomaly time. Such annotations and tool-invocation
success are different from interventional causal truth. The paper's
local timing setup and manufacturing deployment discussion likewise do
not establish banking runtime or treatment-effect performance. We can
adapt its separation of semantic services, deterministic tools and
workflow state as a design pattern while retaining those scientific
limits. \citep{exec13_pulse}

These studies suggest complementary roles: retrieve definitions, draft
alternatives, check logical consistency, construct typed calls, retain
execution artifacts and explain formal results. The adverse findings
point to the same missing link: a linguistic interpretation must be
checked against the intended population, intervention, measurement model
and information clock. An additional critic can help find mistakes, but
agreement between related models is not independent validation of an
assumption.

\hypertarget{a-worked-failure-an-overlap-remedy-changes-the-population}{%
\subsection{A worked failure: an overlap remedy changes the
population}\label{a-worked-failure-an-overlap-remedy-changes-the-population}}

Suppose the bank initially asks for the average invitation effect among
all eligible customers. A diagnostic finds weak overlap for one customer
stratum. An agent proposes trimming that stratum and continues to report
the result as if it answered the original question. Even with perfect
estimation, that continuation can be wrong.

To isolate the target issue, stipulate a three-stratum population with
known masses, propensities and conditional effects. These are generated
numerical assumptions; the units of the conditional effect are thousands
of currency units per customer.

\begin{longtable}[]{@{}>{\raggedright\arraybackslash}p{\dimexpr 0.15\linewidth-2\tabcolsep\relax}>{\raggedright\arraybackslash}p{\dimexpr 0.5\linewidth-2\tabcolsep\relax}>{\raggedright\arraybackslash}p{\dimexpr 0.35\linewidth-2\tabcolsep\relax}@{}}
\toprule
\begin{minipage}[b]{\linewidth}\raggedright
Stratum \(w\)\strut
\end{minipage} & \begin{minipage}[b]{\linewidth}\raggedright
Population mass \(p(w)\), invitation propensity \(e(w)\)\strut
\end{minipage} & \begin{minipage}[b]{\linewidth}\raggedright
Conditional effect \(\tau(w)\)\strut
\end{minipage}\tabularnewline
\midrule
\endhead
\begin{minipage}[t]{\linewidth}\raggedright
0\strut
\end{minipage} & \begin{minipage}[t]{\linewidth}\raggedright
\(0.5,\ 0.05\)\strut
\end{minipage} & \begin{minipage}[t]{\linewidth}\raggedright
\(0\)\strut
\end{minipage}\tabularnewline
\begin{minipage}[t]{\linewidth}\raggedright
1\strut
\end{minipage} & \begin{minipage}[t]{\linewidth}\raggedright
\(0.3,\ 0.50\)\strut
\end{minipage} & \begin{minipage}[t]{\linewidth}\raggedright
\(2\)\strut
\end{minipage}\tabularnewline
\begin{minipage}[t]{\linewidth}\raggedright
2\strut
\end{minipage} & \begin{minipage}[t]{\linewidth}\raggedright
\(0.2,\ 0.80\)\strut
\end{minipage} & \begin{minipage}[t]{\linewidth}\raggedright
\(8\)\strut
\end{minipage}\tabularnewline
\bottomrule
\end{longtable}

The full-population ATE is \(0.5(0)+0.3(2)+0.2(8)=2.2\). Trimming to
true propensities in \([0.1,0.9]\) retains half the population and gives
\(2.2/0.5=4.4\). The conventional overlap-weighted target instead uses
the measure proportional to \(p(w)e(w)\{1-e(w)\}\):

\[
 \psi_{\rm OW}
 =\frac{\sum_w p(w)e(w)\{1-e(w)\}\tau(w)}
        {\sum_w p(w)e(w)\{1-e(w)\}}
 =\frac{0.406}{0.13075}
 =\frac{1624}{523}\simeq 3.1052.
\]

The effect among treated customers is different again:
\(\psi_{\rm ATT}=316/67\simeq4.7164\). Exact rational calculations
reproduce all four values. None is an empirical estimate from a package
or a live language model. Their disagreement arises entirely from the
population measure. \citep{exec13_formulation_targets}

This distinction is visible in a real proposed reuse path. The inspected
oci-agent notebook applies a propensity-based keep mask and averages
AIPW pseudo-outcomes over that subset while labeling the output ``ATO.''
The accompanying paper describes the trimming remedy in the same
terminology. A hard-trimmed average and an \(e(W)\{1-e(W)\}\)-weighted
overlap effect are not generally equal. The integration must either
label the result as a trimmed-population ATE, with its actual
eligibility rule, or implement and validate the intended weighting
operation. It must also retain the original all-eligible request and the
excluded population. Estimated propensities add randomness to selection;
the simple known-law calculation above does not validate that further
inference step. \citep{v12_la_package_oci_agent,v12_la_2607_22443}

The correct translation response is therefore a revised analysis
proposal: the original target is unresolved under the current
information or diagnostics; a specified restricted-population target may
be feasible; and the bank must assess whether that restricted target
still informs its decision. In a public or generated prototype, this
branch can be evaluated without bank data by stipulating both
populations and checking that the system preserves their difference. It
is not necessary to obtain expensive authentic evidence merely to detect
a target substitution.

\hypertarget{concrete-package-choices-and-the-formal-baseline}{%
\subsection{Concrete package choices and the formal
baseline}\label{concrete-package-choices-and-the-formal-baseline}}

The inspected pywhy-llm path is a candidate for eliciting hypotheses and
candidate variables. Its
\texttt{IdentificationSuggester.suggest\_backdoor} delegates to a
language-model confounder suggester. That class's front-door method is a
stub, whereas \texttt{SimpleIdentificationSuggester} has a front-door
method that generates text suggestions. Neither fact establishes a
formal identification algorithm. The proposal should reuse the
suggestion pattern selectively and send the resulting explicit graph and
query to a separately admitted causal engine. It should preserve cited
source passages and unresolved alternatives instead of interpreting a
generated factor list as an approved adjustment set.
\citep{v12_la_package_pywhy_llm}

For oci-agent, the useful candidate is the separation of a
specification, a deterministic analysis notebook and a critic report.
The selected source has a general notebook runner and a target-label
issue requiring correction. A bounded banking interface should expose
fixed reviewed operations with resource limits and typed outputs. The
proposal does not adopt a general notebook executor as the causal MCP
surface. Dependency resolution, failure handling and executable parity
remain separate implementation tasks; the source's research status is
not a reason to infer a deployed bank service.
\citep{v12_la_package_oci_agent}

CausalDS is a candidate for generated scene and task design, but its
grader needs a semantic review before direct adoption. The inspected
identification/scoring helpers distinguish several genuine states, yet
some missing metadata defaults and finite-interval checks require
attention. Unknown identification due to a missing backend is not
scientific nonidentification. A malformed answer is not a justified
abstention. A capped score can conceal severe tail losses unless the
uncapped losses and coverage are also reported. The proposal therefore
takes the task separation and private-reference design while requiring
explicit metadata and failure validation.
\citep{v12_la_package_causalds}

The deterministic formulation baseline has three layers. A schema check
verifies that population, unit, intervention, comparator, horizon, clock
and scale are present. A semantic check compares these fields with the
supplied brief and dictionary: a post-contact duration cannot become a
pre-contact covariate, an invitation cannot silently become actual
uptake, and a requested all-eligible effect cannot become a
recipients-only effect. A formal check evaluates identification under
the explicitly supplied candidate graph and returns an expression,
model-relative nonidentification, unsupported scope or an unresolved
input. These checks validate different things. A well-formed schema does
not prove the story is true, and a correctly evaluated graph does not
prove its edges.

For a generated benchmark, the reference can enumerate a finite set of
admissible formulations \(\mathcal Q(b)\) for brief \(b\). An explicitly
resolved brief may have one acceptable target. An ambiguous brief may
have several and require the system to request or retain the missing
distinction. A baseline can return those alternatives directly from the
supplied reference fields and use exact calculations such as the example
above. This establishes a transparent lower-cost comparator and exposes
what a language model is expected to add: interpretation of unstructured
language, evidence retrieval and useful clarification, rather than
numerical work already handled by a deterministic engine.

\hypertarget{evaluation-before-a-costly-bank-study}{%
\subsection{Evaluation before a costly bank
study}\label{evaluation-before-a-costly-bank-study}}

The initial evaluation can use public data dictionaries and openly
documented financial examples alongside generated causal mechanisms.
Public data support questions about available variables and clocks; they
usually do not supply intervention truth. Generated mechanisms can
supply that truth for the declared model. A reference author should
create matched examples that vary one decisive feature: invitation
versus uptake, all eligible versus treated population, pre-decision
versus post-decision measurement, complete variables versus noisy
proxies, or a single action versus repeated treatment. The target-change
example supplies one exact reference family. This is a proposed
evaluation design; it is not a reported live-agent comparison.

Train/development and final evaluation should be split by the underlying
scenario or structural model, not merely by paraphrases of the same
question. Otherwise a system can learn the reference target from exposed
near-duplicates. All requests should remain in the evaluation
denominator, including ambiguous requests, timeouts, unsupported
estimands, malformed answers and justified nonidentification. Report
target and clock fidelity, unjustified resolution of ambiguity, false
causal answers, valid abstention, numerical error where a target is
actually identified, interval availability and coverage where ground
truth permits them, and total resource use. A high completion rate alone
rewards an agent for confidently answering questions that should remain
unresolved.

Retrieval, structured prompting and an actor/critic loop should be
compared against the same deterministic baseline and fixed task
families. A retrieval ablation tests whether source passages improve
interpretation. A critic ablation tests whether additional review
catches errors, while reporting the added calls, time and newly
introduced mistakes. A graph-perturbation comparison tests whether a
system notices that an assumption changed. Repeated runs reveal
stability. These studies need their actual model, tool, prompt, data and
stopping-rule identities frozen before final evaluation. Agreement
across repeated runs is useful evidence about stability; it is not
evidence that the agreed answer is causally correct.

Several gaps remain scientific. A general method for recovering
stakeholder intent from incomplete language is not established. The
immediate remedy is to preserve admissible alternatives and make the
missing decision explicit. A general guarantee that retrieval or debate
validates causal edges is also absent; the current studies justify
investigating those tools with factual and graph-specific references,
not treating consensus as truth. Identification from noisy proxies,
dependent observations or repeated interventions requires the
corresponding statistical model and method. Finally, whether an assisted
workflow saves bank analyst time or changes an investment decision
requires a later workflow study after a bounded prototype shows merit.
Those gaps narrow the claims and motivate specific experiments; they do
not prevent writing or testing the proposal's public/generated
formulation contract. The next two studies develop decisions within that
contract: which population or weight class the result protects, and how a
clarification workflow resolves a symbolic goal before selecting tools.
Neither statistical robustness nor goal completion can replace an explicit
choice of the intended causal question. \label{exec13:formulation-end}


% SOURCE docs/release/2026-10-29/source/docs/writing-repair/2026-10-29/src/feda358a5f61c756.tex CHAR 0
\section{Whose effect is the assistant being asked to estimate?}
\label{exec22:targets-scope}

A banking assistant translating an offer proposal must specify whose outcome
the offer is intended to improve. ``The average effect is positive'' leaves
open whether the average describes existing customers, a new acquisition
population, an equal mix of branches, or a precision-weighted statistical
summary. These choices can lead to different actions even when every group
effect is known. Vilfort's \emph{Robust Inference for Weighted Estimands}
provides a useful statistical treatment of this ambiguity
\cite{exec22_weighted_vilfort}. It is an econometric contribution to specifying
and checking targets, rather than an algorithm for language-model discovery.
Its value for a causal assistant begins after the group effects have a
defensible interpretation.

\subsection{Separate identification, aggregation and sampling uncertainty}

Let $\theta=(\theta_1,\ldots,\theta_K)'$ contain effects for $K$ customer groups,
with each entry measured in the same units and over the same horizon. A weight
vector $w$ defines
\begin{equation}
 \tau_w=w'\theta,\qquad \mathbf{1}'w=1.
 \label{exec22:targets-definition}
\end{equation}
Normalization alone permits negative entries. If the question concerns an
actual customer mixture, require the simplex
$\Delta_K=\{w:w\geq0,\mathbf{1}'w=1\}$. A signed regression contrast may be a
legitimate statistical target, but it should not be described as the mean
outcome of an unspecified population. Group definitions, treatment regimes,
calendar times and horizons must be compatible before aggregation. A weight
change cannot repair an effect whose identification assumptions have failed.

The exact experiment studied in the paper assumes
\begin{equation}
 \widehat\theta\sim N(\theta,\Sigma),\qquad \Sigma\succ0,
 \label{exec22:targets-experiment}
\end{equation}
where $K$ is fixed, $\Sigma$ is known, and the baseline weights and alternative
class are fixed and known. The full joint covariance matters; separate group
standard errors do not supply it. Let $\Lambda$ be a nonempty compact convex
class of alternative normalized weights. It represents stated ambiguity or
disagreement about the target, rather than an automatically identified set of
possible bank populations. A target outside this class receives no protection
from robustness calculated inside it. Choosing the class after inspecting
estimated effects also changes the statistical problem.

Consider an authored two-group illustration with stipulated effects
$\theta=(6,-2)'$ USD per customer and baseline
$w=(3/4,1/4)'$. Its contribution is $4$. The alternative population
$\lambda=(1/4,3/4)'$ has contribution $0$. At a stipulated cost of USD~1 per
customer, the respective net values are $3$ and $-1$. Nothing about this
disagreement requires a hallucination or an inaccurate group-effect estimate.
The populations answer different questions. These are exact project examples,
not observations from a bank or replications of the paper's applications
\cite{exec22_weighted_reference}.

One transparent ambiguity class moves a fraction $\epsilon$ of baseline
weight to any simplex distribution:
\begin{equation}
 \Lambda_\epsilon=(1-\epsilon)w+\epsilon\Delta_2,
 \qquad 0\leq\epsilon\leq1.
 \label{exec22:targets-class}
\end{equation}
Because the smallest and largest group contributions are $-2$ and $6$, its
effect and net-value ranges are
\begin{equation}
 \{\lambda'\theta:\lambda\in\Lambda_\epsilon\}
   =[4-6\epsilon,4+2\epsilon],\qquad
 \mathcal V_\epsilon=[3-6\epsilon,3+2\epsilon].
 \label{exec22:targets-value}
\end{equation}
The worst-case net value is strictly positive only for $\epsilon<1/2$.
The alternative population above first enters the class at $\epsilon=2/3$:
its first coordinate requires
$1/4\geq(1-\epsilon)3/4$. The break-even ambiguity level therefore has a clear
business meaning. It states how much customer-mix uncertainty the proposed
decision tolerates under stipulated effects. It does not estimate how uncertain
the true mix actually is.

Even a correctly calibrated interval for the baseline need not protect the
other population. Under the additional stipulated model
$\widehat\theta\sim N((6,-2)',I_2)$, the baseline estimator has variance $5/8$.
Its ordinary equal-tail 95\% interval is
$w'\widehat\theta\pm z_{.975}\sqrt{5/8}$. Coverage of the alternative target
$0$ is the analytic probability
\begin{equation}
 \Phi\!\left(z_{.975}-\frac{4}{\sqrt{5/8}}\right)
 -\Phi\!\left(-z_{.975}-\frac{4}{\sqrt{5/8}}\right)
 \simeq0.0009686482.
 \label{exec22:targets-wrong}
\end{equation}
Here $\Phi$ is the standard normal distribution function. Thus the interval
retains 95\% coverage for its declared baseline while covering this other
target about 0.0969\% of the time. This is an exact known-normal calculation,
not a simulation, a defect in baseline inference, or an execution of the
paper's robust procedure. It illustrates why an assistant must preserve the
population as part of the result.

\subsection{Use heterogeneity to understand weight sensitivity}

The paper separates effect heterogeneity from the distance between targets.
Define the covariance-metric heterogeneity and distance by
\begin{equation}
 H(\theta)=\left\{\min_\gamma
  (\theta-\mathbf{1}\gamma)'\Sigma^{-1}
  (\theta-\mathbf{1}\gamma)\right\}^{1/2},\qquad
 d_\Sigma(\lambda,w)=\{(\lambda-w)'\Sigma(\lambda-w)\}^{1/2}.
 \label{exec22:targets-geometry}
\end{equation}
The minimum centers the effects at their generalized least-squares constant.
Since $(\lambda-w)'\mathbf{1}=0$, subtracting any constant leaves the target
difference unchanged. Cauchy--Schwarz in the whitened coordinates gives the
paper's sharp bound
\begin{equation}
 |\tau_\lambda-\tau_w|
 \leq H(\theta)d_\Sigma(\lambda,w).
 \label{exec22:targets-bound}
\end{equation}
A banking interpretation follows directly: populations can disagree little
because effects are nearly homogeneous, because their weights are close in the
relevant metric, or both. None of these statements checks confounding.
In the example with $\Sigma=I_2$, the centering constant is $2$,
$H^2=32$, and $d_\Sigma^2=1/2$. The squared bound is $16$, exactly the squared
effect difference $(0-4)^2$. This attained witness checks the stated geometry.

Minimax-bias weights minimize the maximum covariance-metric distance to
$\Lambda$. For the two-group class at $\epsilon=2/3$, the first coordinate
ranges from $1/4$ to $11/12$. With identity covariance the center is
$7/12$, with effect $8/3$ and worst-case effect deviation $8/3$. This is the
center of the stipulated oracle range $[0,16/3]$. It illustrates the geometric
compromise, but minimizing this bias radius need not minimize total interval
length, and the compromise weights need not be the population the sponsor
wishes to serve. A precision or robustness preference is a decision input;
an assistant should expose it rather than silently substitute it for the
business target.

\subsection{State which confidence event is protected}

The paper constructs a one-sided upper confidence bound $U_\beta$ on
$H(\theta)$, with error budget $\beta$, from the generalized least-squares
residual statistic. Under the exact experiment its square has a noncentral
chi-square law with $K-1$ degrees of freedom and conventional noncentrality
$H(\theta)^2$. The paper's notation indexes this distribution by
heterogeneity itself. A software conversion must reconcile that notation with
the noncentrality accepted by a numerical library. Confusing $H$ and $H^2$
would change the bound. The source also specifies a zero-bound branch when
the observed statistic is sufficiently small; arbitrary root finding without
that branch would not implement the stated procedure.

Let $R_w=\max_{\lambda\in\Lambda}d_\Sigma(\lambda,w)$,
$B_\beta=U_\beta R_w$ and $s_w=\sqrt{w'\Sigma w}$. The paper's two-sided
interval uses the folded-normal critical value $c_{1-\alpha}(b)$, defined as
the $(1-\alpha)$ quantile of $|N(b,1)|$:
\begin{equation}
 C_w=\left[ w'\widehat\theta
       \mathbin{\pm}c_{1-\alpha}(B_\beta/s_w)s_w\right].
 \label{exec22:targets-ci}
\end{equation}
This combines sampling variation with a confidence bound on target
disagreement. The heterogeneity bound, class optimization and critical-value
calculation each require their own numerical checks.

Proposition~6 of the selected original establishes
\begin{equation}
 \Pr_\theta\{\lambda'\theta\in C_w\}\geq1-\alpha-\beta
 \quad\text{for every fixed }\lambda\in\Lambda\text{ and every }\theta.
 \label{exec22:targets-point}
\end{equation}
This uniform statement about each fixed target differs from the event that
one interval contains every target together. Appendix~C.1, Proposition~7,
gives the two-sided simultaneous guarantee
\begin{equation}
 \Pr_\theta\{\lambda'\theta\in C_w\ \text{for all }\lambda\in\Lambda\}
    \geq1-2\alpha-\beta.
 \label{exec22:targets-set}
\end{equation}
The one-sided simultaneous guarantee is $1-\alpha-\beta$. With
$\alpha=\beta=.05$, the stated lower guarantees are 90\% for each fixed
target and 85\% for the simultaneous two-sided event. They are lower bounds,
not assertions that actual coverage equals those numbers. To allocate a 5\%
budget to the latter event one possible choice is
$\alpha=.0125$, $\beta=.025$. Reporting ``95\% confidence'' without stating
the event and both budgets would conceal a substantive change in protection.

This distinction matters for agents that inspect many customer mixes before
selecting a favorable one. A guarantee for every preassigned target does not
automatically guarantee coverage of a target selected from the same results.
Simultaneous protection can protect such a selection within a fixed admissible
class, under its assumptions and its actual confidence budget. It does not
license learning the class from outcomes without further justification, nor
does it correct other forms of model or treatment selection.

\subsection{Move to feasible estimation without discarding the conditions}

The source develops asymptotic extensions with a uniform central limit theorem,
consistent covariance and weight estimates, eigenvalue bounds and regularity
for estimated classes. These conditions are material for dependent banking
panels, where covariance estimation is part of the scientific argument.
The main asymptotic development keeps $K$ fixed. A growing-$K$ appendix
requires separate technical assessment before use here. Treating a large
collection of sparse branch effects as though it automatically satisfied the
fixed-dimensional experiment would bypass that issue.

Estimated classes introduce another boundary. Consistency of a class does not
ensure that it contains every fixed target on the limiting boundary in a given
sample. The source's Proposition~U8 uses a containment condition; an enlarged
class with fixed positive slack is a sufficient route in its example.
Remark~8 explicitly distinguishes the theorem for any fixed $\delta>0$ from
applications reported at $\delta=0$. The observation that $\delta=.0001$
leaves displayed intervals unchanged is a numerical rounding observation,
rather than a proof of the theorem at zero slack. A proposed implementation
must state which class, containment argument and slack it actually uses.

For event studies, the required inputs are the group-time effect vector and
its full joint covariance, with compatible aggregation definitions. An interval
for each event time does not automatically become simultaneous protection
across all event times. The paper's public-data applications illustrate its
authors' results; without acquiring and inspecting their data, code and terms,
they do not establish a replication or a ready software bridge. A bounded
comparison should first use public known-law group effects, verify the oracle
geometry, exercise covariance and class-boundary failures, and then compare
the feasible procedures on matched targets.

The proposal should offer three explicit choices: commit to a substantively
justified fixed population; report a sensitivity range over declared weights;
or add a justified robust confidence construction when sampling uncertainty
also matters. The library, CLI and MCP representations must carry the group
definitions, weight role, full covariance, class, target-selection rule,
confidence event and budgets. A language model can help elicit these fields
and explain their consequences. It cannot infer missing population truth from
well-written prose or turn a valid interval for one target into evidence for
another.


% SOURCE docs/release/2026-10-29/source/docs/writing-repair/2026-10-29/src/2fd495e1a3825bde.tex CHAR 0
\section{Goal clarification before causal analysis and tool dispatch}
\label{exec23:gist-scope}
A request to improve a distressed-loan portfolio does not yet specify whether
the decision is to change contact frequency, offer a payment holiday, alter
collections or merely explain observed arrears. Dispatching a scientifically
valid estimator for the wrong intervention still answers the wrong business
question. GIST-CMTF addresses a related upstream problem: inferring a symbolic
goal before filtering the tools exposed to a language-model agent. Its June
2026 original develops candidate goals, ambiguity detection and clarification
within a common precondition--effect vocabulary. The authentic ten-page
original and a cutoff-selected repository revision are useful design evidence,
but their workflow use of causal does not by itself identify a treatment effect
or supply a structural causal model for borrowers.
\citep{exec23_gist_paper,exec23_gist_sources}

\subsection{What the state-transition model supplies}
\label{exec23:gist-model}
Represent a tool by $t_i=(d_i,R_i,E_i,c_i,\rho_i)$: a description, required
state predicates, produced predicates, optional cost and optional risk label.
A state $s$ and goal $g$ are sets of predicates. A symbolic tool is executable
when $R_i\subseteq s$, and its declared transition adds $E_i$ to the state.
The goal-inference layer maps the user's request and state to candidate goals
and confidence scores; the paper describes confidence and margin rules that
can request clarification instead of committing to one goal. Clarification is
itself represented as a transition, for example from an ambiguous request to
a specified goal. This supplies a concrete place to represent missing business
variables in a CLI or MCP request.
\citep{exec23_gist_paper}

The proposed banking extension should require an intervention, comparator,
population, outcome, horizon, information clock and decision owner before
constructing an effect query. A lender's goal predicate such as
\emph{servicing option selected} does not establish exchangeability, support
or a defensible adjustment set. Those are separate scientific obligations.
Likewise, a workflow edge saying a tool produces a graph does not establish
that the graph is causally correct. The useful inheritance is an explicit
contract for ambiguity and clarification, with scientific identification
review following goal selection.
\citep{exec23_gist_paper,exec23_gist_sources}

The inspected runner implements breadth-first search over monotone symbolic
states and returns the first tool on a shortest path, with ties determined by
tool order. Optional cost and risk labels in the conceptual representation
are not thereby optimized. The semantic-goal comparator uses lexical overlap
with candidate descriptions. The model inference prompt requests one selected
goal and its confidence, while the actual clarification rule combines that
output with a fixed confidence threshold, vague-word detection and the
benchmark's annotated request type. Thus the released implementation is a
specific hybrid rule, not an evaluation of an arbitrary calibrated goal
posterior or the complete candidate-margin formulation.
\citep{exec23_gist_sources}

\subsection{Reported results and reproducible aggregation}
\label{exec23:gist-results}
The paper evaluates 120 constructed tasks: 40 explicit goals, 40 ambiguous
goals, 30 missing-variable requests and 10 clarification-required requests.
Six methods and seven backends produce 5,040 saved rows, with deterministic
mocked tool effects. Independent data-only checks confirm 5,040 unique
model--method--task keys, 32 tools, positive saved token totals and the
published aggregate tables. GIST-CMTF succeeds in 815 of 840 rows, compared
with 673 for top-goal CMTF and 696 for the semantic-goal comparator. Its
21 wrong-goal outcomes contrast with 163 under top-goal CMTF. These are
reaggregations of the released final dataset, not independently repeated
model calls or evidence of general superiority on banking requests.
\citep{exec23_gist_paper,exec23_gist_sources,exec23_gist_reference}

The disaggregated results matter. On explicit-goal tasks the paper reports
94.6\% success for GIST-CMTF versus 98.9\% for the semantic-goal comparator,
with unnecessary clarification on 16.8\% of explicit requests. Its reported
advantage is concentrated in ambiguous requests. The final saved dataset
retains 170 labeled errors: 128 precondition failures, 24 tool-parse errors,
14 empty responses, three goal-parse errors and one invalid selection.
The authors state that infrastructure-corrupted rows were cleaned before the
final analysis. The final positive-token file cannot establish the complete
attempt, retry or provider-failure denominator without those original logs
and cleaning records. The paper explicitly leaves latency, monetary cost,
user satisfaction and real multi-turn clarification unmeasured.
\citep{exec23_gist_paper,exec23_gist_sources,exec23_gist_reference}

\subsection{A source-to-result reproducibility gap}
\label{exec23:gist-provenance}
The repository revision dated 30 June 2026 is later than the 15 June paper
but before the literature cutoff. Its identity is therefore distinct from
proof that it generated the published results. Eleven authenticated artifacts
include the runner, task generator, registry, task file, final metrics,
summary tables, citation metadata, dependency declaration and MIT notice.
The requirements are unpinned, and no raw provider trace appears in the
inspected complete tree. The paper's documented replacement of a Nova preview
backend with Nova Pro is visible in the saved rows, while the runner's default
list still names the preview profile. An explicit override could explain the
model list; it must be retained in a reproducible execution record.
\citep{exec23_gist_sources}

Three further discrepancies prevent direct replay claims. First, the runner
forces clarification for each nonexplicit request type and then assigns the
annotated intended goal. Yet 53 of 560 nonexplicit GIST rows in the saved data
report zero clarification. Second, the runner reads a \texttt{gold\_path}
field to count wrong tools. None of the 120 task objects has that field;
all instead provide \texttt{gold\_chains}. With those unmodified inputs,
that branch has no expected next tool against which to increment its count,
whereas the saved file contains 3,520 wrong-tool counts. Third, the runner's
written CSV columns differ from the released metrics columns. An unreleased
producer, transformation or earlier revision could reconcile these artifacts;
the available bytes do not establish which. This is an explicit provenance
gap, not an allegation that the reported observations were fabricated.
\citep{exec23_gist_sources,exec23_gist_reference}

Benchmark-assisted clarification also limits transport to real users. The
runner consults annotated request type and resolves a clarification by
substituting the intended goal, without observing a possibly ambiguous user
answer. The all-tools and state-aware comparators receive the intended goal
as prompt context. A production assessment must separately test learned
ambiguity detection, actual user responses and the absence of such labels at
inference time. Open-world goals, imperfect state predicates and failed or
nondeterministic tools remain explicit limitations in the original paper.
These limitations preserve the value of the design hypothesis while preventing
a claim that workflow goal validation solves general causal reasoning.
\citep{exec23_gist_paper,exec23_gist_sources}

\subsection{A public-evidence qualification route}
\label{exec23:gist-decision}
Retain GIST-CMTF as a goal-contract and evaluation reference. Before proposing
direct integration, reconcile the exact generating revision, task schema,
metric definitions, model overrides and cleaning log. A project-owned
prototype can then compare a label-free goal classifier, a clarification
policy, a simple lexical baseline and an explicitly named oracle on held-out
public or generated business requests. Freeze the scenario definitions before
outcome inspection; retain wrong goals, unsupported causal questions,
nonresponses, retries and abstentions in the full denominator. Actual
clarification responses and their costs must be represented, rather than
silently replacing them with the correct goal. This protocol is prospective;
no model is invoked or downloaded code executed for this source review.
\citep{exec23_gist_paper,exec23_gist_sources}

An illustrative cost equation makes the missing economics concrete. Let
$\Delta p$ be a descriptive difference in success frequency, $\Delta T$
the additional tokens and $\Delta C$ the additional clarifications per task.
For stipulated success value $v$, token price $c_T$ and clarification cost
$c_C$, define $\Delta N=v\Delta p-c_T\Delta T-c_C\Delta C$.
The saved GIST versus top-goal averages give
$\Delta p=71/420$, $\Delta T=45979/420$ and $\Delta C=23/70$.
With $v=10$ dollars and $c_T=1/1000$ dollar, clarification cost USD 2 gives
$\Delta N=388021/420000$ dollars, whereas cost USD 6 gives
$-163979/420000$ dollars. These exact stipulated amounts illustrate a decision
threshold; they do not estimate bank ROI or causal performance gains.
\citep{exec23_gist_reference}

The scientific owner remains responsible for translating a clarified business
objective into a causal estimand and diagram. The maintainer owns schema
compatibility, raw traces and versioned provenance; the operator owns the
actual endpoint, authentication and measured resource costs. A verified
workflow predicate should record what evidence established it. It must not
transform a language model's declaration that a goal or graph is correct into
an accepted scientific premise. The route is sufficiently concrete to justify
a bounded prototype proposal, while native capability, independent acceptance
and broader frontier coverage remain open.

\subsection{A resolved goal still needs a causal target}

The chapter's examples establish a common requirement at different levels.
An incentive slope is not a finite offer contrast; a trimmed or robustness
weighted average is not automatically the owner's population; a literature
link is not an identified effect; and a completed symbolic goal is not a
scientific conclusion. Each substitution can change the action even if the
underlying calculation or tool sequence is correct.

Useful assistance records those distinctions, exposes what is unresolved and
preserves the evidence behind each accepted commitment. The subsequent formal
analysis then addresses the declared target under stated assumptions. The
next chapter applies this discipline to deposit pricing, arrears servicing,
underwriting and branch campaigns before developing their measurement and
monetary consequences.

